\section{Related Work}
\label{sec:related}

\subsection{Multi-Member Districts and Nested Structures}

Multi-member district (MMD) systems are well-studied in political science and
electoral law.  \citet{grofman1993} reviews proportional and semi-proportional
MMD systems and their effects on minority representation; \citet{bowler1999}
analyzes strategic voting under MMD arrangements.  Neither line of work
addresses the computational question of constructing hierarchically consistent
maps across chambers.

Nested electoral structures---where one legislative body's districts are unions
of another's---appear in several European parliaments and in some US state
constitutions.  The classic example is the German Bundestag, which pairs
single-member constituencies with proportional list seats drawn from the same
geographic regions.  \citet{taagepera1989} analyzes the seat-product model for
mixed-member systems; \citet{cox2006} examines how nesting affects strategic
entry.  These works treat nesting as a political science fact, not a
computational design problem.

The US context is different: the three chambers are separately apportioned,
separately litigated, and separately optimized.  No prior computational
redistricting literature, to our knowledge, treats consistency across chambers
as an algorithmic objective.

\subsection{Computational Redistricting}

The algorithmic redistricting literature largely focuses on single-chamber
problems.  \citet{altman2011} and \citet{fifield2020} describe Markov chain
Monte Carlo (MCMC) sampling over single-chamber district plans; \citet{chen2013}
uses simulated annealing to generate unbiased district plans for gerrymandering
detection.  Graph-partitioning approaches based on METIS \citep{karypis1998}
are widely used in practice; they optimize edge-cut as a proxy for compactness.

Recursive bisection methods \citep{simon1991} produce hierarchical district
trees naturally, but prior work applies them to a single chamber at a time.
The GeoSection algorithm \citep{B8} extends recursive bisection with
ratio-optimal METIS partitioning; ApportionRegions \citep{B11} provides
integer seat allocation for spine regions.  NestSection builds directly on
both.

The GerryChain project \citep{deford2019} and related ensemble methods focus
on detecting gerrymandering by sampling the space of valid plans.  NestSection
is complementary: rather than detecting gerrymandering post hoc, it constrains
the feasible plan space by design.  The compatibility score $\sigma$ determines
how tightly the constraint binds.

\subsection{Cross-Chamber Gerrymandering}

\citet{gelman1994} note that partisan effects compound across legislative
chambers when maps are drawn by the same partisan majority.  \citet{chen2020}
documents cases where congressional and state legislative maps were drawn
sequentially to reinforce each other's partisan effects in North Carolina,
Wisconsin, and Pennsylvania.

The legal literature on cross-chamber coherence is sparse.  \textit{Rucho v.\
Common Cause}, 588 U.S.\ 684 (2019), held that federal courts cannot police
partisan gerrymandering in congressional maps, but left state courts free to
act under state constitutional provisions \citep{miller2019}.  Several state
supreme courts have struck down single-chamber maps (Pennsylvania \textit{LWV
v.\ Commonwealth} (2018); North Carolina \textit{Harper v.\ Hall} (2022);
New York \textit{Harkenrider v.\ Hochul} (2022)) without addressing the
cross-chamber consistency problem.

\textit{Allen v.\ Milligan}, 599 U.S.\ 1 (2023), reaffirmed the applicability
of Section~2 of the Voting Rights Act (VRA) to single-member congressional
districts, but again in a single-chamber context.  \textit{Callais v.\ Landry}
(W.D.\ La.\ 2024) addressed VRA compliance in a state legislative map; the
court's focus on proportionality at a single level highlights the gap that
NestSection addresses: if VRA compliance is evaluated chamber by chamber, a
partisan actor can satisfy each chamber's requirements independently while
maintaining coherent cross-chamber advantages.

\subsection{Factorization and Apportionment Theory}

The connection between prime factorization and hierarchical bisection trees
appears in \citet{B11} (ApportionRegions) and implicitly in the METIS
partitioning literature.  \citet{huntington1921} and \citet{balinski1982}
develop the theory of divisor methods for apportionment; neither addresses
multi-chamber compatibility.  The NestSection compatibility score
$\sigma(C, S, H) = (1 - g/m) \times 100$ is, to our knowledge, a novel
metric for cross-chamber structural alignment.

Divisibility conditions similar to our strict-nesting requirement appear in
local government law: many state constitutions require county-level districts
to be composed of whole precinct units, and precinct units to be composed of
whole census blocks.  NestSection extends this hierarchy upward to the
legislative level.
